\documentclass{article}
\usepackage[utf8]{inputenc}
\usepackage[spanish]{babel}
\usepackage{amsmath} % Para comandos de matemáticas

\title{Evaluación de Electrotecnia - Ley de Ohm y Potencia (Tema 1)}
\author{Prof. Gustavo}
\date{\today}

\begin{document}

\maketitle

\section*{Instrucciones}

\begin{itemize}
    \item Resuelva cada problema mostrando claramente la fórmula utilizada, la sustitución de valores y el resultado final con unidades.
    \item Responda a la pregunta teórica con claridad y concisión.
\end{itemize}

\begin{enumerate}
    \item \textbf{Problema de Potencia ($P$):}
    Una \textbf{resistencia calefactora} de laboratorio es alimentada con una corriente de $I = 4 \text{ A}$. Si la resistencia del elemento es de $R = 25 \, \Omega$. Calcule la \textbf{Potencia disipada} ($P$) en Watts.
    \[ \text{Fórmula Clave: } P = I^2 \cdot R \]
    
    \item \textbf{Problema de Corriente ($I$):}
    Un \textbf{cargador de batería} de celular opera a una tensión de $V = 5 \text{ V}$. Si la resistencia interna del circuito es de $R = 10 \, \Omega$. Determine la \textbf{corriente} ($I$) consumida en Amperes.
    \[ \text{Fórmula Clave: } I = \frac{V}{R} \]
    
    \item \textbf{Problema de Resistencia ($R$):}
    Un \textbf{ventilador de techo} antiguo está conectado a una tensión de $V = 110 \text{ V}$ y se mide una corriente de $I = 5 \text{ A}$. Determine el valor de su \textbf{Resistencia Eléctrica} ($R$) en Ohms.
    \[ \text{Fórmula Clave: } R = \frac{V}{I} \]
    
    \item \textbf{Problema de Tensión ($V$):}
    Un \textbf{soldador eléctrico} industrial consume una corriente de $I = 15 \text{ A}$ y tiene una potencia de $P = 3300 \text{ W}$. Calcule la \textbf{Tensión} ($V$) de alimentación en Volts.
    \[ \text{Fórmula Clave: } V = \frac{P}{I} \]
    
    \item \textbf{Pregunta Teórica:}
    Explique la relación entre la Potencia ($P$) y la Resistencia ($R$) en un circuito, asumiendo que la corriente ($I$) es constante. Mencione el nombre del fenómeno físico que causa la disipación de energía en forma de calor en un conductor.
\end{enumerate}

\end{document}